% sec_introduction.tex -- Introduction
\section{Introduction}\label{sec:introduction}

The Cayley transform
\[
\iota(x)=\frac{1+\ii x}{1-\ii x}
\]
identifies the one-point compactification $\overline\RR$ of the line with
the circle~$\TT$.  Under this identification the normalized Haar measure
on~$\TT$ pulls back to the probability measure
\[
\frac{1}{\pi(1+x^{2})}\dd x,
\]
namely the Cauchy law.  This elementary observation is the starting point
of the paper.  On the geometric side it singles out the Cayley chart among
normalized homeomorphisms of $\overline\RR$ onto~$\TT$.  On the analytic
side it identifies the Haar pullback density with the Poisson kernel at
time~$1$, and hence places the compactification problem inside the
Poisson semigroup on~$\RR$.

The purpose of the paper is to develop the analytic consequences of this
identification.  We first prove that the pullback property rigidifies the
chart: among normalized orientation-preserving
$C^{1}$-homeomorphisms $\overline\RR\to\TT$, the Cayley transform is the
unique one whose pullback of Haar measure is the Cauchy law
(Theorem~\ref{thm:haar-pullback-uniqueness}).  The associated logarithmic
Jacobian yields an entropy identity in Cayley coordinates
(Theorem~\ref{thm:precision-ledger}).  Since the same density is the
Poisson kernel~$P_1$, the problem then passes naturally to Poisson
smoothing on~$\RR$.

The paper has two complementary cores.  The first is algebraic.  Once the
Haar pullback identifies the visible phase factor with a circle, it is
natural to count independent circle factors on the dual side.  We
therefore define, for a finitely generated abelian group~$G$, the circle
dimension
\[
\cdim(G):=\dim\bigl((\widehat G)^0\bigr),
\]
and prove that it is uniquely characterized by natural axioms
(Theorem~\ref{thm:cdim-axiomatic-rigidity}), additive on short exact
sequences (Theorem~\ref{thm:cdim-short-exact-additivity}), multiplicative
under tensor product and $\operatorname{Hom}$, and invisible to
$\operatorname{Ext}^{1}$ (Proposition~\ref{prop:cdim-tensor-hom-ext-laws}).
This turns the circle dimension into a closed algebraic calculus rather
than a mere definition.

The same point of view leads to a finite phase-sampling formalism.  For
$N\ge 1$ we consider
\[
\Sigma_G(N):=\bigl|\operatorname{Hom}(G,\mu_N)\bigr|,
\]
the number of $N$-torsion phase channels visible to~$G$.  We show that
the growth of $\Sigma_G(N)$ recovers the circle dimension
(Theorem~\ref{thm:cdim-phase-spectrum-limit}) and that the full phase
spectrum $N\mapsto \Sigma_G(N)$ determines the group up to isomorphism
(Theorem~\ref{thm:cdim-phase-spectrum-reconstruction}).  In addition, the
Cayley compression gives an operational interpretation of the half-circle
dimension of~$\NN$: dyadic prefix readout at depth~$b$ remains injective up
to scale $2^{b/2}$ and no further
(Theorem~\ref{thm:cdim-operational-half-circle-dimension-N} and
Corollary~\ref{cor:cdim-operational-half-circle-dimension-Nd}).

The second core is analytic.  Its functional-analytic center is a mode
decomposition for the comparison kernel
\[
R_\varepsilon(y):=\frac{1+y^2}{1+(y-\varepsilon)^2}.
\]
After passing to the Cayley angle variable $y=\tan\theta$, the Taylor
coefficients of $R_\varepsilon$ are described by Chebyshev polynomials of
the second kind; see Theorem~\ref{thm:cayley-chebyshev-mode}.  This yields
a finite Fourier expansion for every mode and, in particular, a parity law
showing that the large-time expansion of
$D_{\mathrm{KL}}(P_t*\nu\|P_t(\cdot-\bar\gamma))$ contains only even powers
of~$t^{-1}$ (Theorem~\ref{thm:poisson-kl-even-orders}).

The first concrete analytic consequence is an explicit eighth-order
asymptotic formula for the relative entropy between a Poisson average and
the translated Poisson kernel with matching mean
(Theorem~\ref{thm:poisson-kl-eighth}).  If $h_t=P_t*\nu$ and $g_t$ denotes
the translate of~$P_t$ centered at the mean of~$\nu$, then
\begin{align*}
D_{\mathrm{KL}}(h_t\|g_t)
=\frac{\sigma^4}{8t^4}
+\frac{\sigma^6+6\mu_3^2-8\sigma^2\mu_4}{64\,t^6} \\
&\quad
+\frac{3\sigma^8-12\sigma^4\mu_4+9\sigma^2\mu_3^2+12\sigma^2\mu_6}{256\,t^8} \\
&\quad
-\frac{30\mu_3\mu_5-20\mu_4^2}{256\,t^8}
+o(t^{-8}),
\end{align*}
where $\sigma^2$ and $\mu_k$ are the centred moments of~$\nu$.  The
corresponding coefficients of order $t^{-6}$ and $t^{-8}$ form the first
two levels of a rigidity ladder (Theorem~\ref{thm:poisson-defect-ladder}),
whose unique zero-defect model is again the symmetric two-point law.

The analytic reconstruction theorem shows that the central Poisson trace
and its central spatial derivative recover the Laplace transform of the
centred characteristic function and therefore determine the centred law
uniquely (Theorem~\ref{thm:poisson-central-two-channel}).  Together with
the mean, these data determine the original measure.  In this form the
result separates the even and odd parts of the centred law into two
explicit observation channels.

The same mode calculus also has a canonical Hilbert completion.  The
secant family generated by the comparison kernel has exact Gram kernel
\[
K(a,b)=\frac{2}{4+(a-b)^2}
\]
and yields, through Theorem~\ref{thm:mode-space-rkhs}, a unitary
identification between the zero-mean mode space and the
reproducing-kernel space generated by~$K$.  A further identification with
the Poisson image space $\sqrt{\pi}\,e^{-|D|}L^2(\RR)$
(Proposition~\ref{thm:gram-space-poisson-image}) gives a concrete Fourier
model and a holomorphic extension to the strip
$\{z\in\CC:\ |\Im z|<1\}$.  Thus the RKHS description is not an
external addendum to the asymptotic theory.  It is the Hilbert
completion of the Cayley--Chebyshev mode calculus itself.

The background ingredients are classical.  We use standard facts about
the Poisson semigroup and Fourier transform
\cite{Folland1995,Grafakos2014,Stein1971}, the Moore--Aronszajn theory of
reproducing kernels \cite{Aronszajn1950}, and the classical Stieltjes
inversion formula \cite[Chapter~XVII]{Feller1971}.  For entropy methods in
the stable setting we refer, for example, to the work of
Bobkov--Chistyakov--G\"otze on convergence to stable laws in relative
entropy and to Toscani's fractional-information approach to entropy
decay for stable laws \cite{BobkovChistyakovGotze2013,Toscani2015,Toscani2016}.
These papers address normalized sums and information dissipation in the
stable central limit theorem.  The present paper is different in scope and
in scale: we do not study a limit theorem for sums, but instead derive a
coefficient-level large-time expansion for a single Poisson semigroup orbit
relative to the translated Poisson profile itself.  In particular, the
mode calculus forces a full odd-order vanishing principle, and the first
two non-trivial coefficients admit a defect decomposition that does not
seem to be visible in the available stable-law entropy literature.

The reconstruction theorem should likewise be compared with the classical
Stieltjes inversion principle.  Classically, one recovers a measure from
boundary information on its Cauchy transform.  Theorem~\ref{thm:poisson-central-two-channel}
shows that in the present centred Poisson setting the relevant law can be
recovered instead from two scalar observation channels on the positive
half-line, namely the central Poisson trace and its central derivative.

The paper is organized as follows.
Section~\ref{sec:preliminaries} fixes notation and conventions.
Section~\ref{sec:cayley-gate} records the basic formulas for the Cayley
transform, and Section~\ref{sec:haar-pullback} proves the rigidity of the
Haar pullback.  Section~\ref{sec:circle-dimension} develops the algebraic
circle-dimension package, the phase-sampling reconstruction theorem, and
the operational half-circle interpretation.  Section~\ref{sec:precision-potential}
contains the entropy identity, the Cayley--Chebyshev mode calculus, the
Poisson asymptotics, the centred reconstruction theorem, and the
reproducing-kernel model.  The appendix records the weighted trigonometric
integrals used in the eighth-order expansion.
